\section{Positive spin model}\label{sec:spin}

We convert the signed lattice Hamiltonian into a positive-exchange model
whose interaction strengths can be encoded by distances. Singlet-pair
mediators implement the sign conversion; fixed planar templates then
make the resulting link lengths independently adjustable. The
same-member versus opposite-member sign mechanism is the isotropic
specialization of the mediators in
Piddock and Montanaro~\cite[Section 2.4 and Lemma 5]{PiddockMontanaro2017}.
We keep the quantitative proof here, including the scalar shifts and
the terms unchanged during subdivision.

\subsection{Energy reduction through singlet pairs}

Write \(E\) for the smallest energy in finite dimensions. The following
elementary second-order estimate is a form of the usual effective
Hamiltonian expansion~\cite[Section 3.2]{BDL2011}; we prove precisely the
bottom-energy statement needed here. In particular, the first-order
block \(PBP\) need not vanish.

\begin{lemma}[Second-order bottom-energy estimate]\label{lem:perturbation}
Let \(A,B\) be self-adjoint operators on a finite-dimensional Hilbert
space. Suppose \(A\) has kernel projection \(P\ne0\),
\(A\ge gQ\), where \(Q=I-P\) and \(g>0\), and
\(\|B\|\le g/4\). Then
\[
 E(A+B)=E\big((PBP-PBA^{-1}BP)|_P\big)+O(\|B\|^3/g^2)
\]
where the inverse vanishes on the kernel.
\end{lemma}
\begin{proof}
Indeed \(|E(A+B)|\le\|B\|\). At that energy the \(Q\)-equation can be solved in terms of the nonzero \(P\)-component of a ground vector. The inverse on \(Q\) used there is \((A|_Q+QBQ|_Q-E(A+B))^{-1}\); it differs from \(A^{-1}|_Q\) by \(O(\|B\|/g^2)\) by the resolvent identity. The \(P\)-equation gives the lower estimate. For the upper estimate take a unit ground vector \(p\) for the displayed compression and use \(p-A^{-1}Bp\) as trial vector. Its energy numerator is the compressed energy plus \(O(\|B\|^3/g^2)\), and the normalization changes this result by at most the same order. The statement includes \(Q=0\).
\end{proof}

First delete all extremely small source weights, below a fixed inverse polynomial in absolute value chosen so the deletion error is at most \(\gamma/100\), by \(\|h^s_{ij}\|=3\). Here and below such cutoffs can be rational. The following operation can be made on a set of edges simultaneously, keeping the other terms unchanged. For each replaced edge \(e=LR\), of weight \(J\ne0\), insert a fresh central pair \(A_e,B_e\), with bond \(\Delta h^s_{A_eB_e}\), and two spokes
\[
    \sqrt\Delta\,v_e(h^s_{LA_e}+h^s_{C_eR}),\qquad
    v_e=\sqrt{2|J|},\qquad
    C_e=\begin{cases}B_e& J>0,\\ A_e&J<0.\end{cases}
\]
Here \(\Delta\) is a common large positive integer. Recall that \(h^s_{ij}=2\operatorname{SWAP}_{ij}-I\) with eigenvalues \(-3,1\) on the singlet and triplets, respectively. The new energy is the previous one minus
\[
       \sum_e (3\Delta+3v_e^2/2)
\]
up to error that can be made at most \(\gamma/100\). To verify this shift, penalize excitations of the fresh pairs by \(A=\Delta\sum_e(h^s_{A_eB_e}+3I)\), with singlet kernel for each pair and gap \(4\Delta\); the rest \(B\) has norm at most \(\operatorname{poly}(n)(1+\sqrt\Delta)\) if the entering model has polynomial size and weights. The first-order compression is the unreplaced part.

Only two actions on the same fresh pair contribute to the inverse term, with denominator \(4\Delta\). Pauli products on one member of that singlet have expectation \(\delta_{\mu\nu}\) (components indexed by \(\mu,\nu\)), and products across the pair have \(-\delta_{\mu\nu}\), by the singlet Pauli relations. Thus the negative inverse term gives \(J h^s_{LR}-(3/2)v_e^2 I\) per edge. The estimate just proved permits a polynomially large \(\Delta\).

Apply this conversion first to all surviving signed edges. For each of its spokes make a path of length nine by two subdivision stages, each using the same operation for positive weights (a path of length three per replaced edge). In the first subdivision replace both spokes, and in the second replace all edges of their resulting paths, not the original central-pair bond. Each stage can use its own polynomial scale as there are only three stages total. Thus in the final positive model \(H_+=\sum K_e h^s_{ij}\) there are, per input edge still present, one central bond and two nine-link paths, from \(L\) to \(A_e\) and from \(C_e\) to \(R\). Let \(m\) be the total number of sites (including original vertices); every actual link weight \(K_e\) and its reciprocal are polynomially bounded. There is a sum of scalar shifts \(C_s\) from the operations, such that
\[
           |E(H_+)-E(H_{\rm in})+C_s|\le 4\gamma/100.       \tag{1}\label{eq:1}
\]
All these intermediate weights and shifts are polynomial-time approximable to inverse polynomial error (the square roots need not be exact rational outputs).

\subsection{Independently adjustable contact lengths}

Equation~\eqref{eq:1} completes the energy conversion. To encode its
positive weights by orbital hopping, we now need a geometry in which
changing one link length does not constrain the others.
Here is a planar contact layout in which all the links have independently adjustable lengths; see Figure~\ref{fig:contact-templates}. Start with original grid spacing 17. Orient the edges in the two positive axial directions, and within each edge gadget use planar coordinates along that direction and in a perpendicular direction, origin at \(L\). In the positive input sign case set \(A_e=(8,0), B_e=(9,0)\); in the negative sign case set \(A_e=(8.5,0), B_e=(8.5,1)\). Recall \(C_e=B_e\) and \(C_e=A_e\), respectively. Each nine-link path joins two on-axis points \(8\) or \(8.5\) apart. Orient it along the edge in order of increasing first coordinate. Make the first two, middle (fifth), and last two steps unit forward steps on that direction, and use two unit forward-sloping steps (steps 3 and 4) with first component \(3/4\) for span \(8\), \(7/8\) for span \(8.5\), balanced by two steps (6 and 7) with transverse components negated.

The link lengths now equal one, and nonlinks have distance at least \(1.4\). Indeed consecutive separations projected along a main path/backbone are at least \(3/4\); nonconsecutive sites on it are separated by at least \(1.5\), and the extra side site for a negative sign is at distance at least \(\sqrt2\) from sites other than its central neighbor (the steps immediately there are straight). These bounds also give separation along each entire axial lattice line across successive gadgets. Parallel lines cannot interfere. Transverse offsets are at most \(2\sqrt{1-(3/4)^2}<1.324\); off-axis internal sites have along-axis projection at least \(2+3/4\) away from bounding original grid coordinates. Against sites of a perpendicular orientation this already gives separation greater than \(1.4\). If both internal sites of perpendicular orientations are on-axis, their coordinate differences in the two directions are at least one each. Original vertices other than a gadget's two endpoints are far from its internal sites by grid spacing. This verifies the separation.

The unit-length layout already separates links from nonlinks. We now
show that its link lengths can vary independently without losing that
separation.

There is a fixed \(\varepsilon\in(0,1/100)\) such that \emph{any}
prescribed link lengths in \([1-\varepsilon,1+\varepsilon]\) can be
implemented with original vertices fixed and nonlinks still at distance
\(>1.2\). To prove this, consider one original-edge gadget. We can first
change the central length by moving \(B_e\), then vary each path's
internal sites while its endpoint velocities are prescribed.

For a path with unit step vectors \(e_1,\ldots,e_9\), a linear
dependence with coefficients \(a_1,\ldots,a_9\) among the rows of its
fixed-end length differential would give, at each internal vertex,
\[
                 a_j e_j-a_{j+1}e_{j+1}=0
                    \qquad (1\le j\le8).
\]
All vectors \(a_j e_j\) would therefore be equal. Since the path has
nonparallel steps, they must all vanish. The differential is thus
surjective, as is the full gadget's link-length differential on its
internal-site variables.

Choose a fixed right inverse for each of the two templates and restrict
the internal-site displacements to its image. In these coordinates the
length increment map is \(q\mapsto q+R(q)\), with \(R(0)=0\) and
\(DR(0)=0\). On a small fixed ball, \(\|DR\|<1/2\); for a
sufficiently small prescribed increment \(d\), the iteration
\(q\mapsto d-R(q)\) is a contraction of that ball. Choose the ball
small enough to preserve the nonlink separation, and then fix
\(\varepsilon\) accordingly. These solves have fixed dimension and
algebraic initial data. Approximating the iterations with guard precision
computes positions to any inverse-polynomial error in polynomial time.
Because the original vertices stay fixed, separate gadgets can be solved
independently; their number does not shrink the allowed length interval.
